\documentclass[11pt]{article}
\usepackage[margin=1in]{geometry}
\usepackage[T1]{fontenc}
\usepackage[utf8]{inputenc}
\usepackage{lmodern}
\usepackage{microtype}
\usepackage{parskip}
\usepackage{hyperref}
\emergencystretch=3em

\title{Materials and Methods for the In Vivo Single-Cell RNA-seq Analysis Workflow}
\author{}
\date{}

\begin{document}
\maketitle

\section*{Materials and Methods}

\paragraph{Single-cell RNA-seq data import and metadata harmonization.}
Filtered feature-barcode matrices were imported from Cell Ranger HDF5 outputs for each sample and converted into Seurat objects without applying additional feature- or cell-level thresholds at object creation. Sample-level metadata were read from the study sample information table and appended to each object. Tumor-derived cells were retained by matching cell barcodes to the external ploidy assignment table, whereas the two cell-line samples were handled as special samples because they were not selected through the same ploidy-derived barcode filter. Cell identifiers were prefixed with their sample identifiers before merging to ensure uniqueness across samples. Sample type was assigned from sample names into 2N-cellline, 4N-cellline, 2N-tumor, and 4N-tumor groups, and dose labels were standardized to 0mg/kg, 30mg/kg, and 120mg/kg where applicable.

\paragraph{Quality control and doublet filtering.}
For tumor samples, cells were retained if they were present in the ploidy-selected barcode list. For the two cell-line samples, quality control was performed using RNA feature counts, UMI counts, and mitochondrial transcript percentage. Cells were retained if they had at least 200 detected features, at least 500 UMIs, no more than 20\% mitochondrial UMIs, and feature and UMI counts below the 99th percentile of the corresponding sample-specific distributions. Doublets in the cell-line samples were identified with scDblFinder using the raw RNA count matrix represented as a SingleCellExperiment object; only cells classified as singlets were retained. Samples with fewer than 100 cells after QC were not processed by scDblFinder and were retained with an explicit low-cell-count skip label.

\paragraph{Data normalization, integration, dimensionality reduction, and clustering.}
Each filtered sample was normalized with SCTransform using the RNA assay, while regressing out mitochondrial transcript percentage and retaining all genes in the SCT assay. Integration features were selected with Seurat's SelectIntegrationFeatures using 3,000 features. SCT integration was performed by preparing objects with PrepSCTIntegration, identifying canonical correlation analysis anchors with FindIntegrationAnchors using \texttt{normalization.method = "SCT"} and \texttt{reduction = "cca"}, and generating an integrated assay with IntegrateData. Principal component analysis was run on the integrated assay with at least 50 PCs computed, and the first 30 PCs were used for neighbor graph construction, clustering, and UMAP embedding. Shared nearest-neighbor clustering was performed with Seurat FindClusters at resolution 0.6, and UMAP coordinates were generated from the same 30-dimensional PCA space.

\paragraph{Composition and UMAP visualization.}
Integrated objects were visualized using UMAP embeddings colored by sample, cluster, sample type, and downstream manual cluster labels. Cluster composition plots were generated as stacked bar charts from Seurat metadata, including sample-type and dose-level summaries where those annotations were available. UMAP plots were generated without rasterization when publication-quality vector output was required, and both PDF and PNG formats were produced for key plots.

\paragraph{Cell-cycle scoring and candidate cell-cycle cluster identification.}
Cell-cycle state was quantified on the RNA assay using Seurat CellCycleScoring and the updated 2019 Seurat S-phase and G2/M-phase gene sets. Gene symbols in the object were matched to the reference cell-cycle gene lists before scoring. For each cluster, the mean S score, mean G2/M score, and fractions of cells assigned to G1, S, and G2/M phases were summarized. Candidate cell-cycle-associated clusters were identified by integrating three criteria: outlying cluster-level cell-cycle scores, enrichment of S-phase or G2/M-phase cell fractions above the across-cluster mean plus one standard deviation, and separation along PCA components whose absolute correlation with S or G2/M scores was at least 0.30. If no PC passed the correlation threshold, the PC with the strongest cell-cycle correlation was used as a fallback. Clusters flagged by at least two of these criteria were considered candidate cell-cycle-related clusters.

\paragraph{Differential expression analysis for experimental comparisons.}
Differential expression analyses comparing sample types, dose groups, and tumor ploidy groups were performed with Seurat FindMarkers on the RNA assay after RNA normalization and scaling where required. Sample-type comparisons included pairwise contrasts between 2N-cellline, 4N-cellline, 2N-tumor, and 4N-tumor groups. Dose comparisons were performed separately within 2N-tumor and 4N-tumor subsets, and 2N-tumor versus 4N-tumor contrasts were also performed within each dose level. Cluster-specific 2N-tumor versus 4N-tumor contrasts were run only when both groups contained at least three cells in the cluster. These comparison-level analyses used \texttt{min.pct = 0.1}, \texttt{logfc.threshold = 0.25}, and the Wilcoxon test. Volcano plots used adjusted \(P < 0.05\) and log fold-change cutoffs of 0.5 for visualization, and the most strongly changed genes were displayed by violin plots and heatmaps.

\paragraph{Cluster marker detection.}
Cluster marker genes were identified by one-versus-rest differential expression using Seurat FindMarkers on the RNA assay data slot. For final reclustered and manually merged cluster sets, the default Seurat FindMarkers test was used, with \texttt{min.pct = 0.1} and \texttt{logfc.threshold = 0} when the goal was to retain a broad marker table for downstream enrichment, and \texttt{logfc.threshold = 0.25} when marker tables were used for focused visualization or merge diagnostics. Differential expression tables were ordered by adjusted \(P\) value and absolute log fold change. Positive markers were defined as genes with adjusted \(P < 0.05\) and positive log fold change, and the top markers per cluster were selected by descending log fold change with adjusted \(P\) value used as a secondary ordering criterion.

\paragraph{Cluster expression similarity analysis.}
Pairwise cluster similarity was quantified from cluster-versus-rest differential expression tables. Gene identifiers were normalized by removing genome-build prefixes and version suffixes, converting to uppercase keys, and collapsing duplicated gene keys by retaining the row with the largest absolute log fold change and then the smallest adjusted \(P\) value. Similarity was evaluated under multiple feature definitions, including all ranked genes, lenient filtering (\(P_{\mathrm{adj}} < 0.05\), \(|\log_2FC| \geq 0.25\), and \(|\Delta \mathrm{pct}| \geq 0.05\)), moderate filtering (\(P_{\mathrm{adj}} < 0.01\), \(|\log_2FC| \geq 0.50\), and \(|\Delta \mathrm{pct}| \geq 0.10\)), strict filtering (\(P_{\mathrm{adj}} < 0.001\), \(|\log_2FC| \geq 1.00\), and \(|\Delta \mathrm{pct}| \geq 0.20\)), and top-\(N\) signed gene sets selected after lenient filtering. The principal top-\(N\) rule retained the top 50 up-regulated and top 50 down-regulated genes ranked by log fold change. Pairwise similarity metrics included directional Jaccard similarity, signed log-fold-change cosine similarity on the union of genes with missing values set to zero, and signed log-fold-change Spearman correlation.

\paragraph{Cluster-level quality control.}
Cluster-level technical quality was evaluated from per-cell RNA counts, detected feature counts, mitochondrial percentage, ribosomal percentage when ribosomal genes were available, log10 genes per UMI, UMI per gene, dominant gene fraction, and the fraction of expression contributed by the top 50 genes. For each metric, cluster-level median values were summarized and compared across clusters using robust z-scores based on the median absolute deviation. Suspicious directions were predefined as low \texttt{nCount\_RNA}, low \texttt{nFeature\_RNA}, low log10 genes per UMI, high mitochondrial percentage, high UMI per gene, high dominant gene fraction, high top-50 gene fraction, and high ribosomal percentage. In parallel, one-versus-rest logistic regression models were fit to assess whether QC metrics alone predicted cluster membership, and McFadden's pseudo-\(R^2\) was reported. Clusters with at least three suspicious QC metrics were labeled high concern, clusters with at least two suspicious metrics or McFadden's pseudo-\(R^2 \geq 0.10\) were labeled moderate concern, and the remaining clusters were labeled low concern.

\paragraph{UMAP-neighborhood-based cluster refinement.}
Potential subclusters 4c and 9c were identified from candidate clusters 4 and 9 using UMAP-space neighborhood information. A reference region was defined from core cell-cycle clusters 6, 10, 11, and 12 together with neighboring clusters 0 and 3. For each candidate cell, the fraction of its 50 nearest UMAP neighbors belonging to the reference region and to the core cell-cycle clusters was computed. Candidate seed cells required at least 60\% reference-region neighbors, at least 20\% core-cell-cycle neighbors, and location inside the convex hull of the reference region. To avoid selecting only boundary cells, connected-component expansion was then performed within each full candidate cluster using a 5-nearest-neighbor graph; cells in the same connected component as selected seed cells were included in the refined 4c or 9c labels.

\paragraph{Manual cluster merging and validation.}
Manual cluster merging was performed after reclustering to consolidate transcriptionally similar clusters. In the final merge step, original reclustered groups 8, 9, and 11 were merged into group 8; groups 1 and 10 into group 1; groups 3, 2, and 7 into group 3; and groups 0 and 12 into group 0. After merging, all resulting groups were renumbered from 0 using consecutive integer labels, yielding nine final clusters. Merge behavior was evaluated by UMAP visualization, one-versus-rest marker detection for the original and merged groups, top-marker overlap, and signed log-fold-change cosine similarity between merged and original clusters. The cosine validation used the same top50-each feature definition as the cluster similarity analysis, retaining up to 50 up-regulated and 50 down-regulated genes after lenient filtering.

\paragraph{Final cell filtering, reintegration, and reclustering.}
After preliminary refinement and merge testing, cells assigned to clusters 9, 4, 3, and 9c were removed before the final clustering analysis. The retained cells were reconstructed as sample-specific Seurat objects from the RNA count matrix and reprocessed from the beginning rather than simply reclustering the previous integrated assay. The retained cells were normalized with SCTransform while regressing mitochondrial percentage, integrated with SCT-CCA using 3,000 integration features, reduced by PCA, clustered with the first 30 PCs at resolution 0.6, and embedded by UMAP. This generated a fresh integrated object for final marker detection, manual merging, and annotation.

\paragraph{Hallmark gene-set enrichment and cluster annotation.}
Cluster annotations were assigned using MSigDB Hallmark gene sets obtained through msigdbr for Homo sapiens. This workflow used over-representation analysis rather than rank-based GSEA. For each final cluster, one-versus-rest marker genes were filtered to adjusted \(P < 0.05\), \(|\log FC| \geq 0.25\), and \(|\Delta \mathrm{pct}| \geq 0.05\), and only the top 100 up-regulated genes ranked by log fold change were used as the query set. The background universe was defined as all genes detected in the RNA assay after symbol cleaning. Hallmark sets with 15--500 genes and at least three overlapping query genes were tested by a one-sided hypergeometric test, and \(P\) values were adjusted using the Benjamini--Hochberg method. Only Hallmark terms with FDR \(< 0.05\) were considered for annotation scoring.

\paragraph{Integrated pathway annotation score.}
Significant Hallmark terms were ranked using an integrated annotation score designed to combine enrichment with marker-level expression evidence. For each enriched pathway, marker strength was measured as the mean absolute log fold change of overlapping marker genes, marker detection was measured as the mean within-cluster detection fraction (\texttt{pct.1}) of overlapping genes, and overlap degree was measured as the mean of overlap/query and overlap/pathway fractions. These three raw metrics were min-max scaled within each cluster and combined with equal weights: one third marker strength, one third marker detection fraction, and one third overlap degree. The highest-ranked Hallmark terms were used as primary and multi-label cluster annotations and were visualized on UMAP, dot plots, and heatmaps.

\paragraph{Caching and reproducibility.}
Differential expression steps were written to reuse existing marker tables when the expected output file was already present, and missing cluster-level tests were run only for clusters lacking cached results. Several cluster-level differential expression steps were parallelized with forked workers on Unix-like systems, with a default request of eight workers where configured. Random seeds were set before stochastic procedures in the main workflows, and session information was saved for the initial integration step. All major intermediate Seurat objects, marker tables, enrichment tables, QC summaries, and figures were written to step-specific result directories to support reproducible downstream review.

\end{document}
